\ssscp{Summary}{14-04-03-05}
Reconstructions for each subgroup are presented in \trf{14-11},
based on data in appendix \ref{ch:SurMar}
and the discussion above.
Unrelated etyma such as **kusay ‘cuscus’,
*kusor ‘cuscus’, and **kVndoR(a) ‘cuscus’ are not addressed here
(see \srf{14-03} for those).

\begin{table}[h]
	\wg\centering\tcp{Etyma for `cuscus' in \tsc{Muna-Buton} and Wallacean subgroups}{14-11}
	\st{6pt}\begin{tabularx}{\linewidth}{l|ll>{\raggedright\arraybackslash}X}\lsptoprule
Subgroup						&	Reconstruction	&	gloss	&	Comments	\\	\midrule
\tsc{Muna-Buton} 		&	**mantəR	&	bear cuscus	&	*əR > **əy > \it{e}	\\
\tsc{Sula-Buru}			&	**mantəR ?	&	cuscus	&	Ambelau could be a loan or retention	\\
\tsc{Ambon-Seram}		&	**mantəR	&	cuscus	&	irregular final V in five languages	\\
\tsc{STB}						&	**mantəR/	&	cuscus	&	**məntəR in Kei and	\\	\hhline{~}
										&	**məntəR	&					&	Fordata only	\\
Banda isolate	&	**məntəR	&	cuscus	&	likely influence from Kei	\\
\tsc{Aru}						&	**məntəR	&	cuscus	&		\\
\tsc{Timor-Babar}		&	**mantəR/	&	cuscus	&	potential irregular final C	\\	\hhline{~}
										&	**məntəR	&	cuscus	&	in \tsc{SW. Maluku}	\\
\tsc{Central Timor}	&	**mantəR	&	cuscus	&		\\
\tsc{Bima-Lembata}	&	\mc{3}{l}{(n./a.; outside range of marsupials)}	\\
\lspbottomrule
	\end{tabularx}
\end{table}

Several things can be seen from these etyma in \trf{14-11}
and the data in appendix \ref{ch:SurMar}:

\begin{itemize}
		\item	The gloss for **mantəR is indisputably ‘cuscus’.
					There is no support for a gloss of ‘bandicoot’ outside of Oceanic languages.
		\item	The final vowel of an umbrella etymon should be *ə.
					Only a limited number of \tsc{Ambon-Seram} languages attest a
					different vowel, of which only one points to *a.
		\item	There is good support for final *R in **mantəR through
					well-attested, regular sound changes.
					The evidence supporting *r in *mans(aə)r is left to guesswork
					since there is so little data supporting the proto-phoneme as
					a whole, particularly in word-final position.
		\item	There is widespread support for medial *nt or *nd in **mantəR.
					If we needed to choose an over-arching form to represent the
					data in the \tsc{Muna-Buton} languages
					and the Wallacean subgroups,
					we would choose **mantəR ‘cuscus’ over **mandəR,
					since the *nd forms are more easily explainable by the common
					phenomenon of post-nasal voicing.
					But we are not positing a proto-form in a common parent language;
					we are merely accounting for the close forms in many subgroups
					in the Wallacean region
					and in southeast Sulawesi reflected in \trf{14-11}.
					In particular, shared irregular penultimate *a > **ə in Aru,
					two neighbouring languages of \tsc{Seram-Tanimbar-Bomberai},
					Banda, and two languages of Timor points to this term spreading
					by contact to some extent.
		\item	The etyma in \trf{14-11} show a remarkable consistency
					in both form and meaning, with only trivial differences.
\end{itemize}